\section{Discussion}
\label{sec:discussion}

% Evidence map: paper/analysis.md (Observed/Inferred/Speculative); paper/conclusion.md;
% canonical macros; Results (Sec.~\ref{sec:results}). No new experiments. No invented citations.

We revisit the research question under the frozen CIFAR-10 protocol of
Section~\ref{sec:setup}:
whether replacing the ResNet-20-A post-activation residual block by a full
pre-activation unit reduces best test error.
Throughout, we separate \emph{observed evidence} (what the recorded runs show),
\emph{interpretation} (what those outcomes suggest in context),
\emph{hypothesis} (plausible mechanisms that remain unproven),
and \emph{unsupported causal claims} (statements not licensed by the present design).

\subsection{What the Baseline Establishes}
\label{sec:disc-baseline}

\paragraph{Observed.}
Under the shared recipe, ResNet-20-A attains best test error $\ErrResNet\%$,
whereas Plain-20 attains $\ErrPlain\%$
(Section~\ref{sec:results-residual}; Table~\ref{tab:main-results}).
The residual model is lower by $\ResidualAdvantagePP$ percentage points
($\DeltaResNetMinusPlain$\,pp for ResNet $-$ Plain).

\paragraph{Interpretation.}
Within this depth-20, single-seed CIFAR-10 setting, the residual shortcut is associated
with a clear reduction in best test error relative to the matched plain network.
This supports treating ResNet-20-A as a usable frozen control for subsequent architecture
substitutions, rather than as an under-implemented baseline.

\paragraph{Unsupported causal claim.}
These two runs do not by themselves prove residual learning is universally beneficial
beyond the tested depth, dataset, and recipe, nor do they identify which shortcut
implementation details are necessary for the observed gap.

\subsection{Why the Pre-activation Replacement Did Not Improve the Baseline}
\label{sec:disc-preact}

\paragraph{Observed.}
Full PreAct-ResNet-20-A reaches $\ErrPreAct\%$ best test error, which is
$\DeltaPreActMinusResNet$ percentage points \emph{higher} than ResNet-20-A
(Section~\ref{sec:results-preact}).
The preset improvement hypothesis required a reduction of at least $0.30$\,pp
relative to the frozen control; that threshold was not met
(Section~\ref{sec:setup}; Phase~3 design record summarized in \texttt{paper/conclusion.md}).

\paragraph{Interpretation.}
The substitution is an architecture change that also relocates BN and ReLU relative to
convolution, while the training recipe (learning rate, milestones, epochs, augmentation,
Option~A, depth) remains the post-activation baseline recipe.
Therefore the negative outcome is evidence about
\textbf{pre-activation as a drop-in replacement under the frozen recipe},
not evidence that pre-activation is intrinsically inferior in all training regimes.

\paragraph{Hypothesis (unproven).}
A plausible account---consistent with, but not proven by, the recorded trajectories---is
that the frozen schedule is poorly matched to the PreAct unit on this shallow network,
so that the intended benefit of cleaner identity mappings does not appear as lower test error
in this single run.
\texttt{[NEEDS\_VERIFICATION]}: controlled learning-rate / schedule retuning for PreAct
was not performed in the present evidence bundle.

\paragraph{Unsupported causal claim.}
We do \emph{not} endorse blanket dismissals of pre-activation residual units
(for example, claims of intrinsic inferiority across recipes, depths, or datasets).
No multi-seed estimate, no PreAct-specific hyperparameter search, and no deeper-network
sweep are available here to support such generalizations.

\subsection{Evidence from the Ablations}
\label{sec:disc-ablation}

\paragraph{Ablation~A --- observed.}
Removing the first $\mathrm{BN}{+}\mathrm{ReLU}$ yields $\ErrAblationA\%$ best test error.
Relative to full PreAct, this is $\DeltaAblationAMinusPreAct$\,pp
(magnitude $\DeltaAblationAVsPreAct$\,pp).
Relative to ResNet-20-A, the difference is $\DeltaAblationAMinusResNet$\,pp
(Section~\ref{sec:results-ablation}).

\paragraph{Ablation~A --- interpretation.}
The large recovery versus full PreAct indicates that, in this frozen setting, the
\emph{presence of the first pre-activation stage} co-occurs with the worse PreAct outcome.
The small gap versus ResNet-20-A is an observed single-seed delta only; it does not license
claims of a statistically validated improvement over the frozen baseline.

\paragraph{Ablation~A --- hypothesis vs.\ unsupported claim.}
\emph{Hypothesis:} the first BN+ReLU may be mismatched to the shallow CIFAR-10 depth-20 network when
trained with the post-activation recipe.
\emph{Unsupported:} we have not identified a causal pathway (gradients, feature statistics,
or optimization geometry) that proves why the first stage hurts; no interventional
experiment beyond this one architectural ablation is recorded.

\paragraph{Ablation~B --- observed.}
Removing the second $\mathrm{BN}{+}\mathrm{ReLU}$ produced numerical instability
(status \AblationBStatus) and no comparable best test error
(Section~\ref{sec:results-ablation}).

\paragraph{Ablation~B --- interpretation.}
Under the current learning rate $0.1$ recipe, this configuration is not a valid accuracy
comparator. The failure constrains what can be said about ``accuracy contribution'' of the
second stage: the present evidence speaks to \emph{trainability under this recipe}, not to a
ranked test-error effect.

\paragraph{Ablation~B --- unsupported claim and future check.}
We do \emph{not} conclude that the second BN is universally indispensable for all depths,
optimizers, or learning rates.
A natural follow-up is a controlled rescue experiment (for example, lower learning rate or
warmup) to test whether Ablation~B can become trainable; such a rescue is
\texttt{[NEEDS\_VERIFICATION]} / future work and is not claimed as completed here.

\subsection{Interpretation of the Negative Result}
\label{sec:disc-negative}

\paragraph{Hypothesis status.}
The improvement hypothesis---that full PreAct would strictly reduce best test error relative
to frozen ResNet-20-A under the shared protocol---is \textbf{not supported} by the recorded
evidence ($\ErrPreAct\%$ vs.\ $\ErrResNet\%$; $+\DeltaPreActMinusResNet$\,pp).

\paragraph{Proper scope of the negative result.}
What is supported is a scoped statement:
\textbf{the hypothesized improvement was not supported under the frozen training recipe and
single-seed setting.}
What is \emph{not} supported is any blanket dismissal of pre-activation residual units
outside that scope.

\subsection{Research Implications}
\label{sec:disc-implications}

\begin{enumerate}
  \item \textbf{Architecture and recipe are coupled.}
  Evaluating a block reordering while freezing the post-activation schedule answers a
  drop-in question; it does not answer the best attainable PreAct accuracy after retuning.
  \item \textbf{Negative results shrink the hypothesis space.}
  Ruling out ``free gain from full PreAct under this recipe'' focuses follow-ups on
  recipe rescue, multi-seed estimation, and finer ablations, rather than on claiming a
  universal architectural ranking.
  \item \textbf{Next experiments should be controlled, not rhetorical.}
  Priority directions already suggested by the evidence bundle include multi-seed repeats
  for ResNet-20-A vs.\ Ablation~A, and LR/schedule rescue attempts for PreAct / Ablation~B
  (\texttt{paper/analysis.md}, \texttt{paper/conclusion.md}).
  These remain proposals until executed and registered in the experimental SoT.
\end{enumerate}
